\section[Conclusion and Outlook]{\hspace{0.4em}Conclusion and Outlook}%
\label{sec:conclusion}

We have developed a first-order Picard iteration for the characteristic
initial value problem for the Einstein vacuum and Einstein--massless-scalar
equations in double-null gauge.  The construction follows the geometric
dependency structure of the equations and combines characteristic constraint
solves with spectral elements in the two null coordinates, pole-free angular
operators, and independent first-order residual audits with separate
metric--connection consistency checks.  This provides a unified
numerical framework for exact benchmarks and for characteristic data with no
explicit interior solution.

The regular exact tests cover vacuum Schwarzschild and Kerr geometries and the
scalar Fisher--JNW family.  Metric errors remain small in regular domains and
across a Schwarzschild horizon.  The Schwarzschild-interior experiment also
identifies the present method's loss of accuracy as the curvature singularity
is approached.  Beyond the explicit families, the vacuum calculations evolve
a strong nonspherical outgoing perturbation and crossed characteristic data
with singular limiting initial curvature.  The Einstein--scalar calculations
further demonstrate
coordinate convergence for nonspherical, low-regularity characteristic data
on the protected numerical region.

The final experiment gives numerical evidence for a trapped region and
reconstructs its apparent horizon.  On the curved characteristic domain, both
angular discretizations identify the section
\begin{equation}
  (u,v)=(-0.4696271025,0.04),
\end{equation}
where the two future null expansions are strictly negative on the entire
sphere.  Eighteen independently reconstructed MOTSs form a sampled
apparent-horizon tube over $0.00441\leq v\leq0.05$.  Coordinate refinement
changes its $v=0.04$ graph by $1.50\times10^{-7}$ in maximum norm, while an
independent higher-angular-band evolution changes its mean location by
$7.80\times10^{-10}$.  The distinction between exact metric errors,
independently evaluated curvature residuals, trapped-section sign tests, and
the graph MOTS equation remains essential when interpreting the experiments.

Three directions are particularly important for future work.  First, the
angular MOTS solve should be extended with adaptive coordinate refinement
toward the singular endpoint and across alternative foliations, thereby
resolving the anisotropic apparent horizon more completely.  Second, varying
the scalar-pulse and anisotropic-perturbation
amplitudes would reveal the threshold for trapping and the dependence of the
first trapped location on the initial data.  Third, improved coordinates,
adaptive refinement, and more robust iteration strategies are needed to
improve the approximation near curvature singularities.  These developments
would extend the present double-null framework from a local numerical
construction toward a more precise study of anisotropic black-hole formation
and singular spacetime geometry.
